\chapter{Krawtchouk Transforms}
\label{ch:14}
\section{Orthogonal transformation associated to $K$}
\label{sec:14-1}
% TODO: Develop this section; add verified references, examples, and exercises.

\section{Exponential function in the Krawtchouk basis}
\label{sec:14-2}
% TODO: Develop this section; add verified references, examples, and exercises.

\section{Krawtchouk transform}
\label{sec:14-3}
% TODO: Develop this section; add verified references, examples, and exercises.

\section{Krawtchouk transform and FFT}
\label{sec:14-4}
% TODO: Develop this section; add verified references, examples, and exercises.

\section{Binomial bases and dual transforms}
\label{sec:14-5}
Let $P_N=(\binom mn)_{n,m=0}^N$, $J$ be coordinate reversal, and $D_N=\operatorname{diag}(2^m)$. For the binomial column $b_m(n)=\binom mn$,
\[
\Phi_Nb_m=2^m b_{N-m},\qquad P_N^{-1}\Phi_NP_N=JD_N.
\]
This is a nonorthogonal basis. Rayleigh quotients in these coordinates must use the Gram matrix $P_N^TP_N$.

